\section{Ranked candidate gaps}
\label{sec:gaps}

\subsection{How a candidate becomes a thesis gap}

An unanswered question is not automatically a useful contribution.  In
this review, a candidate is retained only when it satisfies four
conditions:

\begin{enumerate}
  \item it is not answered by the supplied corpus;
  \item CRESCO8 provides an unusual or well-controlled setting in which
        to answer it;
  \item the causal variables and outcome can be measured with available
        tools; and
  \item a result in either direction would still be informative.
\end{enumerate}

The ranking below is deliberately conservative.  ``Not answered by the
corpus'' is the current evidence statement.  ``Novel in the field''
remains to be established by a broader literature search before the
thesis proposal is frozen.

\subsection{Priority 1: controlled memory-tier evidence}
\label{sec:gap-primary}

\paragraph{\researchq{1}.}
\emph{For quantized CPU-only LLM inference on one Xeon Max~9480 socket,
how do verified local-HBM and local-DDR placement affect prefill and
sustained decode under otherwise identical conditions?}

This is the primary gap.  Na et al.\ measure LLM inference on Xeon Max
but do not isolate memory tier or quantized \lcpp{} execution.  Ibeid et
al.\ isolate memory behaviour but not LLM inference.  Goto et al.\ stage
selected HPL data in HBM but do not report an isolated HBM effect.
The experiment therefore connects evidence that is currently separate.

The comparison has two pre-registered hypotheses:

\begin{itemize}
  \item \(H_{1a}\): when the measured live set fits in local HBM,
        sustained decode obtains a larger relative improvement than
        prefill because its low-batch matrix--vector work has lower
        arithmetic intensity.
  \item \(H_{1b}\): the observed performance ratio is smaller than the
        local STREAM bandwidth ratio because not all inference time is
        spent streaming the relevant data.
\end{itemize}

Both hypotheses may be rejected without invalidating the study.  A
negligible HBM effect accompanied by low memory bandwidth would show
that the chosen phase is not bandwidth-bound.  A negligible effect
despite high bandwidth pressure would direct attention to cache reuse,
kernel coverage, synchronization or placement errors.  A result close
to the STREAM ratio would be surprising and would require especially
careful validation.

The minimum sufficient contribution is a phase-resolved comparison with
raw data, uncertainty, live-set accounting, page residency and achieved
bandwidth.  It requires no runtime patch and should be completed before
any secondary gap.

\subsection{Priority 2: the capacity and KV boundary}
\label{sec:gap-capacity}

\paragraph{\researchq{2}.}
\emph{How does the HBM benefit change as the verified live set
approaches and crosses usable local-HBM capacity, and which component
of the live set causes the transition?}

This question distinguishes at least three regimes:

\begin{enumerate}
  \item weights, KV state and transient allocations fit in local HBM;
  \item weights fit but context, batch or packed buffers push the full
        live set beyond the tier; and
  \item weights alone exceed usable local HBM.
\end{enumerate}

The relevant independent variables are model/quantization, prompt
length, generated length, batch or concurrent sequences, and KV
format.  The dependent variables remain TTFT, ITL, phase throughput and
achieved bandwidth, with peak resident memory and tier-resolved pages
as mechanism checks.

The main threat is mistaking an allocation failure, paging event or
direct-mapped Cache-mode conflict for a smooth capacity effect.
Flat-mode runs should therefore establish the primary boundary.
Cache mode is a separate policy experiment, not a replacement for
verified Flat placement.

This question can be answered in a useful first form with whole-process
placement: compare all-HBM runs that fit against matched DDR runs and
document where the HBM run ceases to be feasible.  Independent weight
and KV placement is needed only for the stronger question of how to
operate beyond that boundary.

\subsection{Priority 3: topology as a controlled factor}
\label{sec:gap-topology}

\paragraph{\researchq{3}.}
\emph{When additional sockets or SNC domains are used, can explicit
execution and page placement avoid the regressions observed with an
unmanaged configuration?}

The corpus motivates this question through a tension: Na et al.\ report
poor naive 96-core and SNC4 inference results, while Ibeid et al.\ and
Goto et al.\ benefit from explicit domain-local mapping.  A useful
experiment should compare named policies rather than repeat a generic
``one versus two socket'' sweep:

\begin{itemize}
  \item one socket and its local memory;
  \item two independent, socket-local replicas, reporting aggregate and
        per-request performance;
  \item one process spanning sockets without deliberate placement, as a
        negative control; and
  \item a partitioned design only if the runtime can express it without
        changing unrelated kernels.
\end{itemize}

SNC4 should be attempted only if the administrators can expose and
support the mode for allocated nodes.  Remote page counts, UPI traffic
and per-domain bandwidth are required to interpret the result.  If
mode changes are unavailable, the thesis should report that operational
constraint and retain socket-local placement as the topology study.

This gap is secondary because a two-socket result is difficult to
explain before the one-socket tier effect is established.

\subsection{Conditional software gap: selective tiering as a test of a
published prediction}
\label{sec:gap-software}

\paragraph{\researchq{4}.}
\emph{On a tier pair where the model weights are not guaranteed to fit
in local HBM, does an explicit CPU-side split of weights and KV state
beat the best unmodified whole-process policy, and how much of the
headroom predicted for heterogeneous-memory KV placement survives on
real hardware without oracle knowledge?}

\paragraph{Why this framing replaces the previous one.}
An earlier version of this section presented independent weight and KV
placement as an open systems idea, motivated by analogy with Aurora's
HPL, whose bulk matrix resides in DDR while selected panel data are
staged into HBM~\cite{goto2026aurorahpl}.  A targeted literature search
in August~2026 shows that the idea itself is no longer open.  Fang et
al.\ formalise dynamic KV-cache placement across a high-bandwidth and a
slower off-package tier as a latency-minimisation problem and derive an
upper bound on what placement can buy~\cite{fang2025kvplacement}.
Adjacent systems work argues the same architecture from other
directions: Zhang et al.\ decouple weight-centric operators from KV
management so that the reusable weight stream stays in a fast tier on
multi-socket CPUs~\cite{zhang2026cacheresident}, and Jang et al.\
exploit the predictable access pattern of weights and prefix state to
place them across a tiered CXL hierarchy~\cite{jang2026itme}.

The thesis therefore cannot claim to propose selective tiering, and
should not try.  What remains open is narrower, and it is testable.

\paragraph{What the existing treatment leaves unanswered.}
\tabref{tab:rq4-delta} states the difference precisely.  The
distinction that matters most is the third row.  Fang et al.\ exclude
weight placement from their optimisation by assumption: weights are read
at every decode step, so they fix the full model in HBM permanently and
treat KV entries as the only scheduling variable.  That assumption is
available to them because their fast tier holds the model with room to
spare.  It is not available here.  A Max~9480 socket exposes 64\,GB of
HBM2e, and only 16\,GB per domain under Flat+SNC4, so a quantized model
of interest may consume most or all of the tier before any KV state
exists.  The case their formulation rules out, namely KV state in HBM with
weights displaced to DDR, is precisely the case CRESCO8 can be made to
exhibit.

\begin{table}[htbp]
\centering
\small
\tabsetup
\caption{What the published treatment assumes and what this thesis can
observe.  The contribution of \researchq{4} lies in the right-hand
column, not in the idea of tiered placement.}
\label{tab:rq4-delta}
\begin{tabularx}{\textwidth}{@{}P{2.55cm}YY@{}}
\toprule
\headrow
\textbf{Dimension} & \textbf{Fang et al.~\cite{fang2025kvplacement}} &
\textbf{This thesis on CRESCO8} \\
\midrule
Evidence
& Behavioural simulator of decode-stage memory accesses, driven by
  attention patterns extracted from real inference; no hardware
  execution.
& Measured execution on one Max~9480 socket, with page residency and
  achieved bandwidth reported per tier. \\
\addlinespace
Phase
& Decode only, stated as an explicit scope limit.
& Prefill and sustained decode resolved separately, as in
  \researchq{1}. \\
\addlinespace
Weight placement
& Fixed in the fast tier by assumption and excluded from the
  optimisation, because weights are read every step.
& An experimental variable.  Weights are not guaranteed to fit, so
  which class occupies the tier is the question rather than a
  precondition. \\
\addlinespace
Tier pair
& GPU HBM3 and CPU LPDDR5X across NVLink-C2C: roughly an order of
  magnitude in bandwidth, crossed by explicit staging.
& On-package HBM2e and DDR5 on one socket: a reported 2.9--3.4\(\times\)
  bandwidth ratio, both tiers directly load/store addressable, placement
  expressed as OS page policy rather than a copy. \\
\addlinespace
Policy class
& Dynamic per-token migration under perfect foresight of future token
  importance, yielding an upper bound of up to \(5.87\times\) over
  static placement.
& Static, auditable, whole-class placement with no oracle.  The
  reported bound is therefore a ceiling on a different machine under an
  unattainable assumption, not a target to reproduce. \\
\bottomrule
\end{tabularx}
\end{table}

\paragraph{Candidate policies.}
The two policies are unchanged, but their status is not.  Transformer
data have different access patterns from HPL panels, so the useful split
must still be measured rather than copied:

\begin{itemize}
  \item weights in HBM and KV/scratch in DDR, prioritising the repeated
        weight stream in low-batch decode.  This matches the assumption
        Fang et al.\ impose, so a result here tests whether that
        assumption is the right one on a CPU tier pair; and
  \item KV state in HBM and weights in DDR when long context makes
        attention-state traffic dominant.  This is the configuration
        their formulation excludes, and it is the one for which no
        measured evidence was found.
\end{itemize}

\paragraph{Entry conditions and what the contribution now is.}
The static audit shows that neither policy is exposed as an ordinary
CPU-side \lcpp{} control at the pinned commit.  Stage~4 therefore still
requires a small runtime change, and it still starts only after
\researchq{1} and \researchq{2} have demonstrated a benefit and a
capacity constraint.  What has changed is the claim attached to it.  The
contribution is no longer the proposal of an abstraction; it is a
measured, hardware-verified test of a placement policy whose value has
so far been established only in simulation, on a tier pair that inverts
the assumption under which that simulation was run.  Whether that test
is the first of its kind is a claim the scope note below declines to
make.  The
implementation remains the smallest auditable control that expresses the
policy, not a general NUMA scheduler.

The implementation succeeds only if it beats the strongest existing
baseline, not merely an unmanaged run.  Required ablations are:
all-DDR, all-HBM where feasible, operating-system interleave or preferred
policy if supported, and each selective placement.  Page residency and
allocation failures must be reported.

\paragraph{What a null result now establishes.}
This reframing improves the position of a negative outcome.  Previously,
finding no benefit from selective placement would have meant only that
an untested idea did not work.  Now it bounds a published prediction:
it would show that on a directly addressable CPU tier pair with a
modest bandwidth ratio, static whole-class placement recovers little of
the headroom that simulation attributes to placement on a wider,
staged, accelerator-side hierarchy, and it would identify which of the
differences in \tabref{tab:rq4-delta} accounts for the shortfall.
That is a reportable result in either direction, which is the condition
this review requires of every retained gap.

\paragraph{Scope note.}
This subsection is the one place where the review departs from its
corpus-bounded rule, because the reframing depends on sources outside
the four supplied papers.  The search behind it was web-first and is
recorded in \texttt{bibliography/README.md}; it under-represents
venue-only publications and must be repeated against the ACM Digital
Library, IEEE~Xplore and the 2026 ISC/SC/IPDPS/ICS proceedings before
the proposal is frozen.

\subsection{Extensions, not core gaps}

\paragraph{Energy.}
Energy per token is useful only if the accessible interface has a clear
scope, resolution and sampling method.  A node-level measurement cannot
be presented as CPU energy without accounting for the rest of the node.
Energy should therefore follow, rather than delay, the validated
performance study.

\paragraph{Online serving and tail latency.}
Arrival processes, scheduling and request interference matter for
deployment, but they introduce a second class of software effects.
They are appropriate after the single-request mechanism and capacity
boundary are understood.

\paragraph{Cache mode.}
Cache mode may provide a no-code overflow policy, but direct-mapped
conflicts and undocumented firmware policy can complicate causal
interpretation.  It is a valuable extension if mode changes are
available, not a prerequisite.

\paragraph{GPU and multi-node comparisons.}
Neither is needed for the memory-tier claim.  A fair GPU serving
comparison would require matched quality, software maturity, batching,
capacity and power boundaries.  Multi-node LLM inference adds network
and parallelisation variables not addressed by the primary question.

\subsection{Priority and dependencies}

\tabref{tab:gap-priority} states the resulting order and the condition
under which each item may begin.  This ordering turns the original
catalogue of many speculative gaps into one primary question, two
evidence-driven follow-ons and one conditional implementation.  It also preserves a complete thesis if
the runtime never needs to be changed.

\begin{table}[htbp]
\centering
\small
\tabsetup
\caption{Recommended priority.  A later item starts only when the
preceding evidence justifies it.}
\label{tab:gap-priority}
\begin{tabularx}{\textwidth}{@{}P{1.25cm}P{3.35cm}P{2.65cm}YY@{}}
\toprule
\headrow
\textbf{Order} & \textbf{Question} & \textbf{Role} &
\textbf{Entry condition} & \textbf{Useful result if null} \\
\midrule
1 & \researchq{1}: local HBM vs local DDR & Core empirical contribution
& Platform and placement validation complete
& Quantifies where HBM does not matter and identifies the limiting
  resource. \\
\addlinespace
2 & \researchq{2}: capacity/KV boundary & Core or secondary empirical contribution
& At least one stable configuration and a measured live-set model
& Establishes the feasible envelope even if no sharp transition is
  observed. \\
\addlinespace
3 & \researchq{3}: topology-aware scaling & Secondary
& Remote-traffic counters and supported topology controls
& Shows whether extra resources are neutral or harmful under explicit
  policies. \\
\addlinespace
4 & \researchq{4}: selective tiering & Conditional empirical test
& Reproducible HBM benefit, capacity pressure, and no existing control
  expressing the policy
& A negative result bounds a published placement prediction on measured
  hardware. \\
\addlinespace
5 & Energy, serving, Cache mode & Extensions
& Validated measurement interfaces and remaining time
& Adds deployment context without being required for thesis completeness. \\
\bottomrule
\end{tabularx}
\end{table}
